\documentclass[11pt]{article}
\usepackage[margin=25mm]{geometry}
\usepackage{fontspec}
\setmainfont{Latin Modern Roman}
\usepackage{amsmath,amssymb,mathtools}
\usepackage{unicode-math}
\setmathfont{Latin Modern Math}
\usepackage{xurl,hyperref}
\hypersetup{colorlinks=true,urlcolor=blue}
\setcounter{secnumdepth}{0}
\setlength{\parindent}{0pt}
\setlength{\parskip}{5pt}
\setlength{\emergencystretch}{3em}
\newcommand{\sref}[2]{\hyperlink{source:#1:#2}{[S#2]}}
\newcommand{\cataloguescope}[1]{\paragraph{Recorded target.}#1}
\newcommand{\R}{\mathbb R}
\newcommand{\C}{\mathbb C}
\newcommand{\N}{\mathbb N}
\newcommand{\Z}{\mathbb Z}
\newcommand{\Q}{\mathbb Q}
\newcommand{\F}{\mathbb F}
\newcommand{\eps}{\varepsilon}
\begin{document}
\section{Coefficientwise indices and a level-two Calabi--Yau/CFT test}\label{statement}

\subsection{Definitions and mathematical statement}

\paragraph{Shared notation.}$\mathbb N=\{1,2,\ldots\}$; the standard number fields have their usual meanings. All additional conventions are stated locally.


\paragraph{Definitions and notation.}An index is a signed difference of kernel dimensions of an appropriate Fredholm operator, or its equivariant/graded analogue. An elliptic genus is a generating series of such indices in compact settings. Its Fourier coefficients may count signed or weighted geometric data. CFT denotes conformal field theory, whose state-space trace supplies a separate physical calculation to compare with the geometric series.

\paragraph{Catalogue wording retained for traceability.}
Is there an index theoretic interpretation of elliptic genus? When it can, how do we compute it and what is the geometric interpretation of the coefficients? Check that this matches CFT calculations.
\paragraph{Scope and source qualification.}All three requests are preserved. The source does not prescribe a particular noncompact extension, target Calabi--Yau, grading or CFT; these must be specified in each claimed comparison. The catalogue title/status was paired with a machine-generated research attempt; that attempt is not peer-reviewed evidence of settlement. The statement covers the original source questions rather than only the title’s subcase.
\subsection{Short English statement}
Interpret and compute elliptic genera as indices, explain their coefficients geometrically, and compare with conformal field theory.
\subsection{Research attempt}
\paragraph{Model, date and investigated scope.}
OpenAI GPT-6 Astra (configured model \texttt{gpt-6-astra}), ultra reasoning; generated 11 October 2026. Derived coefficientwise integrality and cancellation algebra under explicit Fredholm/finite-grading assumptions. No target-specific genus, coefficient geometry or CFT comparison is computed.

\paragraph{Sources and literature check.}
Read the supplied complete definition, catalogue wording, scope qualifications and source list for this canonical ID. The original-source citations are inherited as supplied references; no fresh primary-source or novelty search was performed for this attempt. No current open-status claim is inferred from this calculation.

\paragraph{Coefficientwise index algebra.}Suppose a specified compact model has a formal graded virtual bundle \(\mathcal E(q,y)=\sum_{n\ge0,m\in\mathbb Z}E_{n,m}q^ny^m\), with finitely many \(m\) at each \(n\), and elliptic Fredholm operators \(D_{n,m}\) twisted by those bundles. Its formal index series is
\[Z(q,y)=\sum_{n,m}\operatorname{ind}(D_{n,m})q^ny^m,\quad
\operatorname{ind}D=\dim\ker D-\dim\ker D^*.
\]
Additivity follows directly for direct sums, because kernels and adjoint kernels split. Thus each coefficient is an integer but is a signed count, not automatically the dimension of a state space or a nonnegative enumerative invariant.
\paragraph{A supersymmetric cancellation test.}In a finite-dimensional graded complex \((C^\bullet,d)\), let \(a_k=\dim C^k\), \(b_k=\dim\operatorname{im}(d:C^{k-1}\to C^k)\), and \(h_k=\dim H^k\). Rank-nullity gives \(a_k=b_k+h_k+b_{k+1}\). Alternating summation cancels the \(b\) terms and proves \(\sum_k(-1)^ka_k=\sum_k(-1)^kh_k\). This supplies the algebraic mechanism by which paired states cancel in a protected trace. It also shows that agreement of an ungraded total count is not the desired graded-index comparison.
\paragraph{Gap.}No concrete target manifold, twisted elliptic operator family or CFT Hilbert space is constructed here. To answer the source one must identify those objects, compute their coefficients geometrically, and prove that the CFT charge gradings and cancellations match the same normalization. In noncompact settings Fredholmness and trace convergence add separate obstructions. The calculation is a known coefficientwise consistency test; it establishes neither an elliptic-genus formula for an unspecified model nor agreement with any particular CFT.

\paragraph{Verification and adversarial audit.}
Checked Fredholm finite-kernel assumptions, virtual-bundle additivity, rank-nullity cancellation and signed versus unsigned counts. This was an LLM adversarial audit, not independent human verification; no theorem prover was used. Prompts A then B applied to this ID: Checked Fredholm finite-kernel assumptions, virtual-bundle additivity, rank-nullity cancellation and signed versus unsigned counts. Scope remains limited: Construct matching geometric operators and a CFT state complex with the same gradings, then prove and evaluate their index identity.

Coordinator checked direct-sum index additivity and rank-nullity cancellation of adjacent cochain boundaries. This additional check was model-assisted, not independent human verification.

\paragraph{Tools, process and contributions.}
Python 3 and local shell used to read supplied files and serialize this manuscript; mathematics derived symbolically in the text, with no numerical experiment or formal proof assistant. Contribution prompt A was followed by prompt B. GPT-6 Astra with ultra reasoning drafted the research and performed its stated self-audit; any additional coordinator contribution is identified in the verification and audit. A GPT-6-based Codex coordinator assembled the manuscript and checked repository integration; its exact variant and reasoning configuration were not exposed. No human mathematical contribution, independent human proof check or formal proof certification is supplied. Complete derivations and the consulted source references are included here; private reasoning transcripts are not published.

\paragraph{Final status and first remaining gap.}
\textbf{UNRESOLVED}. Construct matching geometric operators and a CFT state complex with the same gradings, then prove and evaluate their index identity.

\section{Completion Estimate}
\noindent\textbf{COMPLETION: 0\%}
The calculation is an elementary benchmark or conditional reduction; it establishes no claimed new portion of the complete source target.

\subsection{Sources}
\begin{itemize}
\item[\hypertarget{source:statement:1}{S1}] ULAM.AI and the UnsolvedMath Contributors, UnsolvedMath, catalogue entry. Catalogue curation under CC BY 4.0; definitions and English exposition added locally.
\url{https://huggingface.co/datasets/ulamai/UnsolvedMath}.
\item[\hypertarget{source:statement:2}{S2}] American Institute of Mathematics, Gromov-Witten invariants and number theory, source item 3.2.
\url{http://aimpl.org/gromwitnumthry/3/}.
\end{itemize}
\label{end:statement}
\end{document}
